\section{The Real and Complex Number Systems}
\subsection*{Exercises}
\addcontentsline{toc}{subsection}{Exercises}
\begin{exercise}
Let $A$ be a nonempty set of real numbers which is bounded below.
Let $-A$ be the set of all numbers $-x$, where $x \in A$.
Prove that \[\inf A = -\sup (-A).\]
\end{exercise}
\begin{proof}
Since $A$ is bounded below, the infimum of $A$ exists in $\R$. 
If $a \in A$ then $\inf A \leq a$. Hence,
\begin{equation}\label{ineq:ex_1_a}
-a \leq -\inf A \text{ for all } a \in A,
\end{equation}
and it follows that $-A$ is bounded above. Thus, $\sup (-A)$ exists in $\R$.
Since $- \inf A$ is an upper bound of $-A$ from \eqref{ineq:ex_1_a}, we have $\sup (-A) \leq -\inf A$, hence 
\begin{equation}\label{ineq:ex_1_b}
\inf A \leq -\sup(-A).
\end{equation}
Since $-a \leq \sup(-A)$ for every $a \in A$, we obtain $-\sup (-A) \leq a$ for all $a \in A$.
Therefore, $-\sup(-A)$ is a lower bound of $A$ and 
\begin{equation}\label{ineq:ex_1_c}
-\sup(-A) \leq  \inf A.
\end{equation}
Combining \eqref{ineq:ex_1_b} and \eqref{ineq:ex_1_c}, we conclude that $\inf A = -\sup (-A)$.
\end{proof}
\begin{exercise}
Fix $b>1$.
\begin{enumerate}
\item[(a)] If $m$, $n$, $p$, $q$ are integers, $n>0$, $q>0$, and $r=m/n=p/q$, prove that 
 \[
 (b^m)^{1/n}=(b^p)^{1/q}.
 \]
 Hence it makes sense to define $b^r = (b^m)^{1/n}$.
 \item[(b)] Prove that $b^{r+s} = b^r b^s$ if $r$ and $s$ are rational.
 \item[(c)] If $x$ is real, define $B(x)$ to be the set of all numbers $b^t$, where $t$ is rational and $t \leq x$. 
Prove that \[b^r = \sup B(r)\] when $r$ is rational.
Hence it makes sense to define \[b^x = \sup B(x)\] for every real $x$.
\item[(d)] Prove that $b^{x+y}=b^xb^y$ for all real $x$ and $y$.
\end{enumerate}
\end{exercise}
\begin{proof}
\,
\begin{enumerate}
\item[(a)]
Let $b^{mq}=b^{np}=x$.
Since $n$ and $q$ are positive integers and $b>0$, by Theorem 1.21, we have 
\[
b^m=x^{1/q} \text{ and } b^p=x^{1/n}.
\]
It follows that
\[
(b^m)^{1/n}=(x^{1/q})^{1/n}=(x^{1/n})^{1/q}=(b^p)^{1/q},
\]
since $[(x^{1/q})^{1/n}]^{nq}=[(x^{1/n})^{1/q}]^{nq}=x$.
\item[(b)]Let $m$, $n$, $p$, $q$ be integers and $n>0$, $q>0$ and let $r=m/n$ and $s=p/q$. 
Then by the Corollary to Theorem 1.21, we obtain
\begin{align*}
b^{r+s}
&=b^{\frac{m}{n}+\frac{p}{q}}
=b^{\frac{mq+np}{nq}}\\
&=\left(b^{mq+np}\right)^{\frac{1}{nq}}\\
&=\left(b^{mq}b^{np}\right)^{\frac{1}{nq}}\\
&=\left(b^{mq}\right)^{\frac{1}{nq}} \left(b^{np}\right)^{\frac{1}{nq}}\\
&=b^{\frac{mq}{nq}} b^{\frac{np}{nq}}
=b^{\frac{m}{n}} b^{\frac{p}{q}}
=b^rb^s.
\end{align*}
\item[(c)]
Note that by Theorem 1.21, $b^r>0$ if $b>1$ and $r$ is rational.

Let $n$ be a positive integer and let $0<a<1$ be real. We first show that $0<a^{1/n}<1$.

Suppose to the contrary that $a^{1/n} \geq 1$. Then we obtain 
\[
1 \leq \left(a^{1/n}\right)^n=a<1,
\]
a contradiction. Therefore, if $0<a<1$ then $0<a^{1/n}<1$, where $n$ is a positive integer.

Suppose $r$ and $s$ are rational. 
Let $m$, $n$, $p$, and $q$ be integers and let $n$ and $q$ be positive.
Then we write $r=m/n$, $s=p/q$.

Suppose $r<s$. Then
\[
r-s =\frac{m}{n}-\frac{p}{q} =\frac{mq-np}{nq} <0,
\]
hence
\[
0 < b^{mq-np} = \frac{1}{b^{np-mq}} < 1
\]
and 
\[
0 < b^{r-s} = b^{\frac{mq-np}{nq}} = \left(b^{mq-np}\right)^{1/nq} < 1.
\]
Thus, if $r<s$ then
\[
b^r = b^{(r-s)+s} = b^{r-s}b^s < b^s.
\]

If $t$ is rational and $t \leq r$ then $b^t \leq b^r$. 
It follows that $B(r)$ is bounded above by $b^r$ and $\sup B(r)$ exists in $\R$.
Since $b^r$ is an upper bound of $B(r)$, 
\[
\sup B(r) \leq b^r.
\]
Since $b^r \in B(r)$,
\[
b^r \leq \sup B(r).
\]

We conclude that $b^r = \sup B(r)$.

\item[(d)]
Let $t_1$ and $t_2$ be rational. 
If $t_1 \leq x$ and $t_2 \leq y$ then $t_1+t_2 \leq x+y$.
Thus, we have
\[
b^{t_1}b^{t_2} = b^{t_1+t_2} \leq \sup B(x+y) =b^{x+y}.
\]

Fix $t_1\leq x$. Then 
\[
b^{t_2} \leq b^{x+y}/b^{t_1} \text{ for all } t_2 \leq y.
\]
Hence, 
\[
b^y = \sup B(y) \leq b^{x+y}/b^{t_1}.
\]
Since the choice of $t_1$ was arbitrary,
\[
b^{t_1} \leq b^{x+y}/b^y \text{ for all } t_1 \leq x,
\]
and it follows that
\[
b^x = \sup B(x) \leq b^{x+y}/b^y.
\]
Therefore, we obtain
\begin{equation} \label{ineq:Ch.1_ex2_(d)_1}
    b^xb^y \leq b^{x+y}.
\end{equation}

Notice that 
\[
0 < b^{t_1} \leq b^x \text{ for every } t_1 \leq x
\]
and 
\[
0 < b^{t_2} \leq b^y \text{ for every } t_2 \leq y \mathrlap{.}
\]
Hence, for every $t_1 \leq x$ and $t_2 \leq y$, 
\[
b^{t_1+t_2} = b^{t_1}b^{t_2} \leq b^xb^y.
\]

Now, to obtain
\begin{equation}\label{ineq:Ch.1_ex2_(d)_2}
b^{x+y} = \sup B(x+y) \leq b^xb^y
\end{equation}
from above inequality, we need the following statement:
\[
b^x=\sup \left\{b^t \mid t \in \Q, t<x\right\}.
\]

Let us prove it. Notice that
\[
b^n-1=(b-1)(b^{n-1}+ \cdots + 1) \geq n(b-1)
\]
for any positive integer $n$. Put $c=b^{1/n}>1$. Then we have
\[
b-1 \geq n(b^{1/n}-1).
\]
Hence, if $t>1$ and $n>(b-1)/(t-1)$ then
\[
\frac{n(b^{1/n}-1)}{t-1} \leq \frac{b-1}{t-1} < n,
\]
and it follows that
\[
b^{1/n}<t.
\]

Note that if $t>1$ then $b^{1/n}<t$ for sufficiently large $n$.

Let $B_0(x)=\left\{b^t \mid t \in \Q, t<x\right\}$.
We need to show that $b^x$ is the supremum of $B_0(x)$. 
It is clear that $b^x$ is an upper bound of $B_0(x)$.
Thus, it remains to prove that if $y<b^x$ then $y$ is not an upper bound.

Suppose $0<y<b^x$. Then $b^x/y >1$ and it follows that
\[
b^{1/n} <b^x/y
\]
for sufficiently large $n$. 
Hence, we have
\[
y<\frac{b^x}{b^{1/n}}<b^x.
\]
By the density of rationals, choose a rational $t$ such that $x-(1/n)<t<x$. 
Since $x<s<t+(1/n)$ for some rational $s$, 
\[
b^x \leq b^s <b^{t+(1/n)} = b^tb^{1/n}
\]
and
\[
y<\frac{b^x}{b^{1/n}}<b^t.
\]
Thus, $y$ is not an upper bound of $B_0(x)$ and therefore $b^x = \sup B_0(x)$.

Now, let $t$ be a rational and let $t<x+y$.
Pick $t_2 \in \Q$ such that $t-x<t_2<y$ and set $t_1=t-t_2$.
Then for every $t \in \Q$ with $t<x+y$, there are rational $t_1<x$ and $t_2<y$ such that $t=t_1+t_2<x+y$, hence $b^t=b^{t_1+t_2}=b^{t_1}b^{t_2} \leq b^xb^y$.
Thus, we obtain \eqref{ineq:Ch.1_ex2_(d)_2}.

Combining \eqref{ineq:Ch.1_ex2_(d)_1} and \eqref{ineq:Ch.1_ex2_(d)_2}, we conclude that
\[
b^{x+y}=b^xb^y.
\]
\end{enumerate}
\end{proof}
\begin{exercise}
Prove that no order can be defined in the complex field that turns it into an ordered field.
\end{exercise}
\begin{proof}
Suppose to the contrary that there exists an order $<$ on $\C$ which turns $\C$ into an ordered field.
Let $\mathbf{0}$ and $\mathbf{1}$ denote the additive and multiplicative identities of the complex field, respectively. 
Then $\mathbf{0}=(0,0)$ and $\mathbf{1}=(1,0)$. 

Note that $i = (0,1) \neq \mathbf{0}$.
If $x \neq \mathbf{0}$ then $x^2 > \mathbf{0}$. 
It follows that
\begin{align*}
    i^2=(0,1)(0,1)=(-1,0)>\mathbf{0}.
\end{align*}
Since $(-1,0)+(1,0)=\mathbf{0}$, $(-1,0)$ is the additive inverse of $(1,0)$. 
Since $\mathbf{1}>\mathbf{0}$ and if $x>\mathbf{0}$ then $-x<\mathbf{0}$ in any ordered field, we obtain
\[
(-1,0)=-(1,0)=-\mathbf{1}<\mathbf{0}.
\]

This contradicts the Trichotomy Property. 
\end{proof}